% Front matter — tese curta, PT-PT.

\frontmatter
\pagestyle{plain}

\maketitlepage

%----------------------------------------------------------------------------------------
%	DECLARAÇÃO DE INTEGRIDADE
%----------------------------------------------------------------------------------------

\begin{soi}

\vspace{1cm}

Declaro que realizei este trabalho académico com integridade.

Não recorri a plágio nem a qualquer forma de utilização indevida ou falsificação de informações ou
resultados ao longo do processo que conduziu à sua elaboração.

Declaro, por isso, que este documento é original e da minha autoria, não tendo sido utilizado
anteriormente para qualquer outro fim.

Declaro ainda que tenho pleno conhecimento do Código de Conduta Ética do P.PORTO.

\vspace{0.5cm}

\textbf{Uso de Inteligência Artificial.} Em conformidade com as orientações do P.PORTO/ISEP, declaro
que foram utilizadas ferramentas de inteligência artificial generativa ao longo do desenvolvimento
deste trabalho, e declaro-o na sua extensão real. Foram usadas para escrever e reestruturar parte
substancial do código do sistema e dos seus testes, para implementar os procedimentos de avaliação,
para redigir e editar prosa neste documento, e para conduzir revisões que encontraram defeitos no
software e no texto.

O que é meu é o trabalho para onde essas ferramentas foram apontadas: a escolha do problema, as
questões de investigação, as restrições que definem o sistema (apenas fontes gratuitas,
explicabilidade primeiro, e a recusa de prever preços), e todas as decisões sobre o que construir,
manter ou descartar. Sempre que uma medição contrariou uma afirmação minha, retirei a afirmação em
vez de a defender. Revi o conteúdo deste documento, e ele, os erros que nele restem, e a sua defesa
são da minha responsabilidade.

\vspace{1cm}

% A data acompanha a da capa (definida em main.tex). Estava em \today e mudava sozinha a cada
% compilação, o que numa DECLARAÇÃO ASSINADA é pior do que na capa.
% ⚠️ FALTA CONFIRMAR com o orientador a redação exata que a MEIA/ISEP exige para esta declaração.
ISEP, Porto, setembro de 2026

\end{soi}

%----------------------------------------------------------------------------------------
%	DEDICATÓRIA
%----------------------------------------------------------------------------------------
\begin{dedicatory}
À minha família.
\end{dedicatory}

%----------------------------------------------------------------------------------------
%	RESUMO
%----------------------------------------------------------------------------------------

\begin{abstract}
Quando uma ação se move de forma abrupta, um investidor particular faz quase sempre as mesmas
três perguntas: \emph{isto é invulgar?}, \emph{foi a empresa ou foi o mercado?}, e \emph{já aconteceu
antes, e o que se seguiu?} As ferramentas gratuitas mostram a percentagem e ficam-se por aí; os
terminais profissionais respondem, mas a um custo fora do alcance de quem investe as suas poupanças.

Esta dissertação apresenta o \textbf{InvestiGator}, um sistema de alertas financeiros construído
sobre fontes de dados gratuitas, que responde às três perguntas e mostra sempre como chegou à
resposta. O trabalho assenta em quatro técnicas: deteção de anomalias por \emph{z}-score deslizante;
decomposição do movimento em mercado, setor e empresa, com encolhimento do beta pela precisão da
estimativa; recuperação de precedentes por semelhança semântica, com medição do desfecho observado;
e um modelo supervisionado de triagem, treinado sobre um conjunto de dados próprio, para decidir
que notícias merecem interromper alguém.

Cada componente foi avaliado contra linhas de base transparentes, e os resultados são reportados tal
como se revelaram. A deteção normalizada pela volatilidade dispara de forma muito mais consistente
entre empresas do que um limiar fixo. A recuperação semântica supera as linhas de base lexical e
triviais. O modelo de triagem, treinado em 79\,753 exemplos, \textbf{não} superou uma linha de base
que usa apenas volatilidade, um resultado negativo que se reporta, e que valida por medição a
opção pela simplicidade que atravessa todo o trabalho.

O sistema recusa, por desenho, prever preços: explica o que já aconteceu, com a evidência ao lado.
\end{abstract}

\begin{abstractotherlanguage}
When a stock moves abruptly, a retail investor tends to ask the same three questions:
\emph{is this unusual?}, \emph{was it the company or the market?}, and \emph{has this happened
before, and what followed?} Free tools show the percentage and stop there; professional terminals do
answer, at a subscription cost beyond someone investing their savings.

This dissertation presents \textbf{InvestiGator}, a financial alerting system built on free data
sources that answers those three questions and always shows how it reached the answer. The work
rests on four techniques: anomaly detection with a rolling \emph{z}-score; decomposition of a move
into market, sector and company, with the beta shrunk by the precision of its estimate; precedent
retrieval by semantic similarity, with measurement of the observed outcome; and a supervised triage
model, trained on a purpose-built dataset, to decide which news deserves interrupting someone.

Each component was evaluated against transparent baselines, and the results are reported exactly as
they fell. Volatility-normalised detection fires far more consistently across companies than a fixed
threshold. Semantic retrieval beats lexical and trivial baselines. The triage model, trained on
79\,753 examples, did \textbf{not} beat a volatility-only baseline, a negative result that is
reported, and that validates by measurement the simplicity-first stance running through the work.

By design, the system refuses to predict prices: it explains what already happened, with the
evidence beside it.
\end{abstractotherlanguage}

%----------------------------------------------------------------------------------------
%	AGRADECIMENTOS
%----------------------------------------------------------------------------------------
\begin{acknowledgements}
Ao meu orientador, Prof.\ Luís Gomes, e ao meu coorientador, Rafael Silva, por todo o apoio ao longo
desta tese. Por trocarem ideias comigo e me ajudarem a decidir a direção do tema; por me mostrarem
exemplos de aplicações que já existem e me manterem a par do que se vai fazendo na área; e por me
guiarem na elaboração deste documento. Muito do que aqui está tomou a forma que tem porque houve
antes uma conversa com eles.

À minha família, por todo o apoio e por acreditarem em mim. Por me incentivarem sempre, e pela
alegria com que foram acompanhando a construção disto. É, para mim, a melhor parte de o ter
feito.

Aos meus colegas da Sistrade, e à Sistrade, pela flexibilidade que me permitiu seguir este mestrado
a par do trabalho, pelas conversas informais sobre a aplicação ao longo destes meses, e por serem
boa companhia todos os dias.
\end{acknowledgements}

%----------------------------------------------------------------------------------------
%	ÍNDICES
%----------------------------------------------------------------------------------------
\tableofcontents
\listoffigures
\listoftables

\printnoidxglossary[type=\acronymtype, title=Lista de Acrónimos, nonumberlist]

\mainmatter
\pagestyle{thesis}
